\title{Parameter-Efficient Orthogonal Finetuning via Factorization}

\begin{abstract}
The widespread adoption of large foundation models has highlighted the substantial computational cost associated with training such models from the ground up. Consequently, developing effective strategies for adapting these powerful architectures to specific downstream applications has gained importance. In this work, we examine Orthogonal Finetuning (OFT), a structured approach to adapting foundation models for target tasks. Although OFT demonstrates promising transferability, it remains parameter-intensive because orthogonal matrices have high dimensionality. To enhance parameter efficiency, we analyze OFT through the lens of information transmission, identifying essential properties conducive to more efficient parameter usage. Drawing inspiration from the Cooley-Tukey fast Fourier transform's capacity for efficient data transformation, we introduce an orthogonal parameterization grounded in butterfly structures. This reparameterization gives rise to a novel, parameter-efficient adaptation technique: Orthogonal Butterfly (BOFT), which extends the OFT framework by encompassing it as a particular instance. We validate the BOFT approach with comprehensive experiments across diverse domains, including adapting large vision transformers, state-of-the-art language models, and text-to-image diffusion models for a variety of downstream tasks in both vision and language settings.
\end{abstract}

\section{Introduction}

State-of-the-art models such as ChatGPT~\cite{brown2020language,chen2021evaluating} and Stable Diffusion~\cite{rombach2022high} have exemplified the exceptional generalization capacities inherent to large foundation models. As the capabilities of these models have expanded, so too has the sheer number of their parameters—GPT-3, for example, encompasses approximately 175B parameters. This trend renders training from initial randomization ever more costly and less accessible to the broader research community. Hence, further advancements depend on the capability to repurpose pretrained foundation models for novel applications without necessitating full retraining. To this end, researchers have explored several efficient adaptation paradigms: \emph{model finetuning}~\cite{hulora2022,zaken2022bitfit,guo2020parameter,raffel2020exploring,qiu2023controlling,zhang2022adaptive,chavan2023one}, wherein the original network remains largely intact and only select parameters are updated; \emph{adapter tuning}~\cite{rebuffi2017learning,houlsby2019parameter,pfeiffer2020adapterfusion,lin2020exploring,he2021towards}, which augments the model with additional trainable modules; and \emph{prompt tuning}~\cite{li2021prefix,lester2021power}, involving prepending trainable tokens to model inputs. Among these alternatives, model finetuning stands out due to its straightforward implementation and, crucially, its avoidance of added inference delay.

The core idea underpinning model finetuning is to maintain a certain degree of similarity between the pretrained and finetuned models, as assessed by specific metrics, thereby retaining the knowledge acquired during pretraining. A common practice in current finetuning techniques is the use of a reduced learning rate, which serves as a heuristic to limit changes and thus keep the two models closely aligned. Recognizing that weight matrices have inherent structure, a more systematic strategy focuses on preserving relational properties within these matrices—specifically, the pairwise angular relationships among neurons. This perspective motivates the development of Orthogonal Finetuning (OFT)~\cite{qiu2023controlling}, a novel approach wherein all neurons within a given layer are jointly transformed using a single orthogonal matrix, ensuring that the angles between pairs of neurons are rigorously maintained during the finetuning phase. While OFT has shown substantial improvements in generalizability and convergence for tasks such as text-to-image diffusion model adaptation~\cite{qiu2023controlling}, it is hindered by the substantial number of trainable parameters required by full-sized orthogonal matrices, especially in high-dimensional settings. To make OFT more parameter-efficient, a block-diagonal parameterization is employed, which substantially reduces the parameter count. Nevertheless, this increase in efficiency introduces certain limitations: the sparsity imposed by the fixed block-diagonal pattern means that orthogonal operations are performed separately within each block. Although this pattern economizes on parameters, it may also embed undesirable inductive biases (such as limiting the expressivity of the model and preventing it from representing general linear transformations), and, crucially, the optimal choice of sparsity structure for such matrices remains an unresolved issue.

A central challenge in this context is devising a method to create a dense orthogonal matrix without incurring the substantial parameter overhead typically required. Constructing a $d$-dimensional dense orthogonal matrix conventionally necessitates $\mathcal{O}(d^2)$ parameters, making such direct parameterizations impractical for large $d$. Our strategy introduces an alternative solution: building dense orthogonal matrices by composing several factorized sparse matrices. The underlying principle is to reduce the number of parameters by leveraging increased computational steps—trading space complexity for computation time. Because our dense orthogonal matrix is realized as a product of sparse matrices, this composition must be repeatedly executed during each training phase. To provide context for this factorization method, we interpret the generation of a dense orthogonal matrix as analogous to an information transmission scenario. In this analogy, synthesizing a dense orthogonal matrix through sparse matrix products corresponds to propagating information through a grid-based graphical structure. Adopting this information transmission perspective allows us to conceptualize various factorization strategies for sparse matrices, effectively constraining the parameter count while still retaining sufficient expressive power to capture dense orthogonal transformations.

To facilitate parameter-efficient design within our information transmission formulation, we draw upon the butterfly structures observed in the Cooley-Tukey fast Fourier transform algorithm~\cite{cooley1965algorithm}, which efficiently enable information transfer between disparate nodes~\cite{klasing1994broadcasting}. The butterfly topology found in the Cooley-Tukey procedure naturally motivates a specific mode of sparse matrix decomposition known as butterfly factorization. Utilizing butterfly factorization, it becomes feasible to generate a $d$-dimensional dense matrix as a product of $\mathcal{O}(\log d)$ sparse matrices, with each factor containing only $\mathcal{O}(d)$ non-zero entries. By constructing each sparse factor as an orthogonal matrix, we obtain a highly efficient orthogonal parameterization requiring just $\mathcal{O}(d\log d)$ parameters, representing a considerable improvement over conventional dense parameterizations. This compact orthogonal parameterization underpins our introduction of a new parameter-efficient finetuning technique, Orthogonal Butterfly~(BOFT). BOFT encompasses OFT as a specific instance and presents a unified and general approach for orthogonal finetuning. Notably, both block-diagonal and butterfly structures exhibit sparsity, a trait that facilitates the reduction of trainable parameter counts by imposing structural sparsity on orthogonal matrices. It is instructive to compare our method with the low-rank decompositions characteristic of LoRA~\cite{hulora2022}; indeed, both low-rank and sparse matrices are fundamental structured matrix families~\cite{candes2011robust} leveraged to enhance parameter efficiency.

In contrast to OFT's use of a block-diagonal structure, which balances expressiveness and structural regularity, BOFT leverages the butterfly architecture to achieve a more continuous transition between full orthogonal matrices (which allow greater expressivity) and the identity matrix (which ensures higher regularity). This facilitates the identification of a more compact hypothesis space within the family of orthogonal matrices for use in downstream applications. Given the extensive role that the butterfly structure plays in efficient linear operators—such as the discrete Fourier transform and the discrete cosine transform—we propose that this structural prior introduces a form of implicit inductive bias that enhances generalization and mitigates overfitting. Our reasoning stems from the butterfly structure's inherent ability to reproduce a wide range of classical linear transforms, whereas the block-diagonal architecture utilized by OFT is incapable of achieving such flexibility. Our main contributions can be summarized as follows:

\begin{itemize}[leftmargin=*,nosep]
    \item We examine parameter efficiency within orthogonal finetuning by introducing an innovative information transmission paradigm. This paradigm reframes the construction of parameter-efficient dense orthogonal matrices as a problem of information propagation across a graph arranged in a grid.
    \item Drawing on the butterfly structure central to the Cooley-Tukey algorithm, we present Orthogonal Butterfly—an approach to orthogonal finetuning that is both parameter-efficient and aligned with our information transmission framework.
    \item We offer several theoretical perspectives explaining why BOFT achieves substantial reductions in trainable parameters, yet preserves desirable properties of expressiveness and training robustness. The butterfly-based matrix factorization further enables an intriguing interpolation of weights (refer to Figure~\ref{fig:boft_interpolation}).
    \item To our knowledge, this is the inaugural effort to utilize orthogonal finetuning for diverse applications beyond the domain of controllable text-to-image synthesis~\cite{qiu2023controlling}, thereby highlighting BOFT’s promise as a general-purpose model adaptation technique. We demonstrate the efficacy of BOFT across various adaptation challenges, encompassing both computer vision and natural language processing tasks, where it surpasses leading contemporary approaches in terms of both parameter efficiency and generalization performance.
\end{itemize}

\textbf{Parameter-efficient finetuning (PEFT)}. As foundation models (\emph{e.g.}, \cite{brown2020language,radford2021learning,rombach2022high,kirillov2023segment}) grow in scale and capability, there has been mounting interest in developing parameter-efficient strategies to adapt these models for specific downstream tasks~\cite{treviso2023efficient,ding2023parameter,lialin2023scaling}. Numerous PEFT approaches have been introduced~\cite{houlsby2019parameter,aghajanyan2020intrinsic,hulora2022,edalati2022krona,wang2022adamix,gheini2021cross,zaken2022bitfit,guo2020parameter,sung2021training,ansell2022composable,lester2021power,li2021prefix,vu2022spot,he2021towards,mao2021unipelt,karimi2021compacter,liu2022few,sung2022lst,chen2022parameter,jia2022visual,chen2022adaptformer,zhang2022neural,jie2023fact,lian2022scaling,luo2023towards}, with reparameterization-based techniques~\cite{aghajanyan2020intrinsic,hulora2022,edalati2022krona} being particularly pertinent to our study. LoRA~\cite{hulora2022}, for instance, enhances the pretrained weights by appending the product of two low-rank matrices, leading to notable improvements in natural language applications. While LoRA maintains a fixed rank for the supplementary matrices across all layers, subsequent research~\cite{zhang2022adaptive,valipour2022dylora,zhang2023increlora} has explored adaptive rank assignments across layers to optimize parameter allocation. The straightforward nature of low-rank reparameterizations has driven widespread adoption~\cite{dettmers2023qlora,zi2023delta,chavan2023one}. Motivated by the role of hyperspherical energy in understanding generalization~\cite{liu2018learning,liu2021orthogonal}, \cite{qiu2023controlling} introduces orthogonal finetuning (OFT), a method which optimizes an orthogonal matrix to linearly transform neurons within a layer. This approach yields more stable training and improved generalization compared to LoRA when fine-tuning text-to-image diffusion models. Nonetheless, OFT often requires more tunable parameters than LoRA, suggesting a need for further improvement in parameter efficiency. Additionally, the potential of OFT for tasks beyond the adaptation of text-to-image diffusion models~\cite{qiu2023controlling} remains largely unexplored. BOFT addresses OFT’s parameter overhead by employing butterfly factorization, thereby enabling the application of orthogonal finetuning to broader adaptation tasks with enhanced efficiency.

\textbf{Butterfly structures}. The radix-2 Cooley-Tukey algorithm achieves efficient computation of the $N$-point discrete Fourier transform by successively decomposing it into two $\frac{N}{2}$-point transforms, naturally giving rise to the butterfly pattern, which can be represented as a product of sparse matrices known as butterfly matrices. These butterfly matrices have served as parameterizations for orthogonal matrices to eliminate pivoting in Gaussian elimination and to increase computational efficiency~\cite{parker1995random}; they have also been leveraged to stabilize recurrent neural network training~\cite{jing2017tunable} and for kernel approximation~\cite{munkhoeva2018quadrature}. Fast linear transformations with butterfly parametrizations are studied in~\cite{dao2019learning,dao2019kaleidoscope}, while efficient neural network training leveraging butterfly matrices has been proposed in~\cite{chen2022pixelated,dao2022monarch}. Further, butterfly structures have facilitated improvements in fast matrix-vector operations~\cite{michielssen1996multilevel,o2010algorithm}, data-sparse matrix approximations~\cite{li2015butterfly}, and network communications~\cite{klasing1994broadcasting,li2003linear,rajasingh2016transmission}. Unlike prior approaches, our work focuses on using butterfly structures specifically to increase the parameter efficiency of OFT.

\section{An Information Transmission View on Orthogonal Finetuning}

We begin by outlining the foundational concepts of OFT. In order to adapt the pretrained weight matrix during finetuning, OFT reformulates the updated weight matrix as a product of a trainable orthogonal matrix and the fixed pretrained weights. This is in contrast to LoRA, which adds a low-rank matrix to the weights; OFT instead leverages a multiplicative orthogonal transformation. LoRA attains parameter efficiency through low-rank decomposition, whereas the standard formulation of OFT employs a block-diagonal orthogonal matrix~\cite{qiu2023controlling}. Parameter efficiency in OFT improves as the diagonal block sizes decrease. Figure~\ref{fig:comp_lora} provides a visual comparison. The rationale behind using orthogonal transformations for weight adaptation is to retain the angular relationships among neurons~\cite{liu2018learning,liu2021orthogonal,qiu2023controlling}, which in turn helps conserve the semantic information from the pretraining phase. More specifically, for a pretrained linear layer $\bm{W}^0\in\mathbb{R}^{d\times n}$, OFT trains an orthogonal matrix $\bm{R}\in\mathbb{R}^{d\times d}$, thereby altering the computation from $\bm{z}=(\bm{W}^0)^\top\bm{x}$ to $\bm{z}=(\bm{R}\bm{W}^0)^\top\bm{x}$, where $\bm{x}\in\mathbb{R}^{d}$ denotes the input and $\bm{z}\in\mathbb{R}^{n}$ is the resulting output. To start with no deviation from the pretrained weights, OFT initializes $\bm{R}$ as the identity matrix. The orthogonality constraint on $\bm{R}$ is maintained throughout training by using Cayley parameterization~\cite{qiu2023controlling,liu2021orthogonal}: $\bm{R}=(\bm{I}+\bm{Q})(\bm{I}-\bm{Q})^{-1}$, where $\bm{Q}$ is skew-symmetric, satisfying $\bm{Q}=-\bm{Q}^\top$. To facilitate parameter efficiency, OFT introduces a block-diagonal form for $\bm{R}$: $\text{diag}(\bm{R}_1,\bm{R}_2,\cdots,\bm{R}_r)$, where each $\bm{R}_i\in\mathbb{R}^{b\times b}$ is an orthogonal block, and $br=d$. Although block-diagonalization reduces parameter count, it introduces the implicit assumption that each neuron's dimensions (i.e., the columns of $\bm{W}^0$) are split into $r$ groups, with each group being independently transformed by a distinct orthogonal matrix. Despite the observed empirical performance of this structure, dividing neuron dimensions by index lacks a principled basis, making densely structured orthogonal matrices potentially more appealing. Thus, a pertinent question emerges: \emph{Is it possible to design a dense orthogonal matrix that does not compromise parameter efficiency?}

To explore this problem, we suggest expressing a dense orthogonal matrix as a sequence of products involving several sparse orthogonal matrices. In order to offer a consistent and accessible framework for examining sparsity structures within orthogonal matrix factorization, we recast the generation of a dense orthogonal matrix in OFT as analogous to an information transmission task. More precisely, forming a dense matrix $\bm{R}\in\mathbb{R}^{d\times d}$ by multiplying $m$ square matrices, such that $\thickmuskip=2mu \medmuskip=2mu \bm{R}=\bm{B}_m\bm{B}_{m-1}\cdots\bm{B}_1$, can be interpreted as transferring information across a grid of $d\times (m+1)$ nodes, as shown in Figure~\ref{fig:info_trans}. The rationale behind adopting the information transmission analogy comes from viewing a dense $d \times d$ matrix as a full set of possible connections from $d$ input nodes to $d$ output nodes. Specifically, for any entry $R_{ij}$ in the matrix $\bm{R}$, a zero indicates a blocked path for information flow from node $j$ to node $i$, whereas a nonzero value signifies possible transmission. Consequently, the representation of the dense matrix $\bm{R}$ via the sequential product $\bm{B}_m\bm{B}_{m-1}\cdots\bm{B}_1$ aligns with an iterative communication process over the graphs each $\bm{B}_i$ induces, for $i=1,\dots,m$. Information moves according to $\bm{B}_1$ at the initial step, passing through to $\bm{B}_n$ at the final stage. As illustrated in Figure~\ref{fig:info_trans}, the factorization $\bm{R}=\bm{B}_5\bm{B}_4\bm{B}_3\bm{B}_2\bm{B}_1$ is used as a concrete case, with both the respective sparsity patterns and the corresponding graph presented. This graph arises from an explicit unfolding of the matrix product, where $\bm{B}_i$ serves as a connectivity map linking the $i$-th layer nodes to the nodes at level $i+1$. Concretely, each element $(j_1,j_2)$ in $\bm{B}_i$ specifies the presence (or absence, when zero) of a directed edge from node $j_2$ at level $i$ to node $j_1$ at level $i+1$. In order for $\bm{B}_5\bm{B}_4\bm{B}_3\bm{B}_2\bm{B}_1$ to be dense, it is required that every source node at the first level be able to propagate information to every target node at the sixth level. Focusing instead on $\bm{R}=\bm{B}_4\bm{B}_3\bm{B}_2\bm{B}_1$, corresponding to transmission between level-one and level-five nodes, we observe an instance where information from node $1$ cannot reach node $3$—meaning the composition $\bm{B}_4\bm{B}_3\bm{B}_2\bm{B}_1$ yields a zero in the $(3,1)$ entry. This observation also holds for other partial products such as $\bm{R}=\bm{B}_3\bm{B}_2\bm{B}_1$ or $\bm{R}=\bm{B}_2\bm{B}_1$: the associated graph structure mirrors the resulting sparsity patterns. More broadly, ensuring that a matrix $\bm{R}\in\mathbb{R}^{d\times d}$ is fully dense through such factorizations entails that the $m$ factors $\bm{B}_m,\cdots,\bm{B}_1$ instantiate a collection of directed edges on a $d\times (m+1)$ grid, each edge connecting a pair of nodes in consecutive layers (i.e., columns), so that there exists a communication route from any first-layer node to each node at level $m+1$.

The perspective of information transmission offers valuable insights into the sparse structures present in OFT factorization matrices. In Figure~\ref{fig:blk_diag}, the block-diagonal form of $\bm{R}$ found in the standard OFT is depicted. This arrangement reduces the number of parameters that require training but is insufficient for representing a fully dense matrix $\bm{R}$. The aim is to generate a dense orthogonal matrix by assembling $m$ sparse orthogonal matrices, achieving this with the minimal number of effective trainable parameters. Within the information transmission framework, two key requirements emerge for this objective: \emph{(i)} \textbf{dense connectivity}, where each node in the initial layer is connected, through at least one pathway, to every node in the final layer, and \emph{(ii)} \textbf{minimal free edges}, which seeks to keep the overall number of edges as low as possible while preserving orthogonality. Importantly, orthogonality enforces particular restrictions on edges between successive levels. Specifically, for every matrix $\bm{B}_i$ to maintain full rank—a prerequisite for orthogonality—there must exist $d$ edges to create a one-to-one correspondence between nodes of the $i$-th level and those of the $(i-1)$-th level. As a result, there are at least $d$ edges required between adjacent levels (for instance, Figure~\ref{fig:info_trans} shows 4 in its illustration), which are essential for orthogonality and thus not counted toward the total edge count, since these elements are not subject to training (e.g., a $d \times d$ orthogonal matrix with $d$ nonzero elements, each constrained to be $\pm 1$). Because orthogonal matrices demand fewer parameters than unrestricted matrices, imposing orthogonality introduces additional dependencies among the edges. As an illustrative case, in a $2\times2$ orthogonal matrix, if a zero occupies position $(1,1)$, another zero must be present at $(2,2)$; thus, removing a single edge can require the removal of another. Consequently, the feasible connections are influenced by orthogonality, which may either enforce the presence or absence of certain edges. The utility of the information transmission approach in OFT is especially pronounced when examining the support pattern of the resulting orthogonal matrix, since our main interest lies in the trainable nonzero entries of $\bm{R}$, and not in their specific magnitudes.

A straightforward dense connection between two levels involves $\mathcal{O}(d^2)$ edges, effectively corresponding to a single dense orthogonal matrix and resulting in $d^2-d$ connections (for instance, with $d=4$, this gives 12 edges). As demonstrated in Figure~\ref{fig:info_trans}, an example of a plausible matrix factorization comprises $10$ edges in total, which is fewer than the number needed for a dense orthogonal matrix. This approach allows us to analyze the parameter efficiency of OFT graphically, facilitating the construction of admissible factorizations under this perspective. We draw on the structure of butterfly graphs~\cite{cooley1965algorithm} from the Cooley-Tukey algorithm, which efficiently interconnect $d$ source nodes and $d$ destination nodes using only $\mathcal{O}(d\log d)$ edges. For instance, the configuration in Figure~\ref{fig:info_trans} employs $10$ edges for full connectivity, yet the butterfly network accomplishes the same with just $8$ edges. In what follows, we explain how leveraging the butterfly structure can enhance parameter efficiency.

\section{Orthogonal Parameterization by Butterfly Factorization}

Initially introduced in the Cooley-Tukey algorithm for the purpose of fast Fourier transforms, the butterfly structure serves to efficiently propagate a localized frequency-domain modification throughout the entire spatial domain, drawing a conceptual analogy with our information transmission challenge—where each node at the initial level must communicate with every node at the final level. Furthermore, the butterfly structure has emerged as a widely adopted topology in computer networks~\cite{solihin2015fundamentals,li2011linear} for facilitating effective information exchange. Let us assume $k\geq 2$ is a power of $2$. We begin by introducing the \emph{butterfly factor} $\bm{B}^F(k)$ defined as

\begin{equation}\label{eq:but_factor}
\begin{aligned}
    \bm{B}^F(k)=\begin{bmatrix}
    \text{diag}(\bm{d}_1) & \text{diag}(\bm{d}_2) \\
    \text{diag}(\bm{d}_3) & \text{diag}(\bm{d}_4)
    \end{bmatrix}\in\mathbb{R}^{k\times k},
\end{aligned}
\end{equation}

where each $\bm{d}_i\in\mathbb{R}^{\frac{k}{2}}$ for $i=1,\ldots,4$. For $d=2^N$, we then construct the $d$-dimensional \emph{butterfly matrix} $\bm{B}(d)\in\mathbb{R}^{d\times d}$ through the following recursive definition:

\begin{equation}
\begin{aligned}
    \!\!\bm{B}(d)=\tilde{\bm{B}}(d,d)\!\cdot\!\begin{bmatrix}
    \bm{B}_1(\frac{d}{2})\!\!\!\!\! & \bm{0} \\
    \bm{0}\!\!\!\!\! &\bm{B}_2(\frac{d}{2})
    \end{bmatrix}= \tilde{\bm{B}}(d,d)\tilde{\bm{B}}(d,\frac{d}{2})\cdots\tilde{\bm{B}}(d,2),
\end{aligned}
\end{equation}

Here, $\bm{B}_1(\frac{d}{2})$ and $\bm{B}_2(\frac{d}{2})$ denote butterfly matrices with dimensionality $\frac{d}{2}$. The \emph{butterfly component} is subsequently introduced as $\thickmuskip=2mu \medmuskip=2mu \tilde{\bm{B}}(d,k)=\text{diag}(\bm{B}^F_1(k),\cdots,\bm{B}^F_{d/k}(k))$, representing a $d\times d$ block-diagonal matrix in which each block is of size $k$, and where each $\bm{B}^F_i(k)$ (for all $i$) corresponds to butterfly factors specified by Equation~\ref{eq:but_factor}. This setup enables us to parameterize orthogonal matrices using butterfly matrices. Ensuring the orthogonality of every factor $\tilde{\bm{B}}(d,k)$ within the butterfly matrix $\bm{B}(d)$ suffices to enforce the overall orthogonality. Considering first the case where the block size is $2$, the block-diagonal matrix $\tilde{\bm{B}}(d,2)$ can be conveniently parameterized as orthogonal, either via the Cayley transform or $2$-dimensional rotations, aligning with prior approaches~\cite{liu2021orthogonal,qiu2023controlling}. The sparsity structure found in any butterfly component is simply a permutation of that found in $\tilde{\bm{B}}(d,2)$; as a result, each butterfly component can be systematically parameterized as an orthogonal matrix. This method yields an efficient parameterization for orthogonal matrices composed from several $2\times2$ orthogonal submatrices. Extending this approach as in \cite{chen2022pixelated}, a block butterfly component $\tilde{\bm{B}}^b(d,k)$ is defined where each entry $\bm{d}_i$ becomes a $b\times b$ matrix. To maintain the orthogonality of $\tilde{\bm{B}}^b(d,2)$, each $2b\times2b$ block is independently parameterized as an orthogonal matrix. For the remaining butterfly components $\tilde{\bm{B}}^b(d,k)$ with $k>2$, their nonzero patterns correspond to blockwise permutations of $\tilde{\bm{B}}^b(d,2)$, and thus, they may likewise be made orthogonal in a similar manner. By assembling these components, the forward propagation in BOFT can be defined as follows:

\begin{equation}
    \begin{aligned}\nonumber
    \bm{z}=\big(\bm{R}(m,b)\cdot\bm{W}^0\big)^\top\bm{x},~~~~\text{s.t.}~\bigg{\{}\bm{R}(m,b)=\prod_{i=1}^m \tilde{\bm{B}}^b_{(d,i)}~~\&~~\underbrace{\big(\tilde{\bm{B}}^b_{(d,j)}\big)^\top\tilde{\bm{B}}^b_{(d,j)}=\tilde{\bm{B}}^b_{(d,j)}\big(\tilde{\bm{B}}^b_{(d,j)}\big)^\top\!=\bm{I}_d}_{\forall j\in [1, m]}\bigg{\}},
    \end{aligned}
\end{equation}

where, for convenience, we define $\tilde{\bm{B}}^b(d,2^{m - i+1})$ as $\tilde{\bm{B}}^b_{(d,i)}$, and let $\bm{I}_d$ represent the $d \times d$ identity matrix. The orthogonal matrix $\bm{R}(m,b)\in\mathbb{R}^{d\times d}$ is constructed as the product of several orthogonal butterfly matrices. For notational simplicity, we refer to BOFT with the orthogonal transformation $\bm{R}(m,\frac{b}{2})$ as BOFT($m$,$b$), provided $b\geq 2$. Specifically, if we set $m=1$, then BOFT($1$,$b$) reduces to the block-diagonal OFT~\cite{qiu2023controlling} where each block is of size $b$. When the block size is set to $b=d$, BOFT($1$,$d$) recovers the original OFT~\cite{qiu2023controlling} that uses a fully dense, unconstrained orthogonal matrix. With the choice $m=\log \frac{2d}{b}$, BOFT($\log \frac{2d}{b}$,$b$) makes use of a block butterfly matrix $\bm{B}^{\frac{b}{2}}(d)$ as its $\bm{R}$ and as such, yields a dense orthogonal transformation. Generally, the number of learnable parameters involved in BOFT($m$,$b$) is $\frac{1}{2}(b-1)dm$ when adapting a linear layer of dimension $d\times n$. In the special case where we use the butterfly structure with $m=\log d, b=2$, BOFT requires only $\mathcal{O}(d\log d)$ parameters. This stands in contrast to the conventional OFT approach that utilizes a fully dense orthogonal matrix, incurring a parameter cost of $\mathcal{O}(d^2)$, while the block-diagonal OFT employing $r$ blocks needs $\mathcal{O}(bd)$ parameters. Consequently, achieving a dense orthogonal mapping with the standard OFT mandates $b=d$, whereas BOFT is more flexible, capable of attaining this property for any $b$.

\textbf{Identity initialization for BOFT}. Standard practice in finetuning dictates that one initialize from the weights of a pretrained network, ensuring the finetuned version does not deviate drastically from its original state. LoRA, as an example, initializes the low-rank adaptations to zero. In the context of BOFT, every butterfly factor is initialized as an identity matrix (that is, the underlying skew-symmetric matrices involved in the Cayley transform are set to zero).

\textbf{Multiplicative Dropout for BOFT}. The LoRA method~\cite{hulora2022} employs a Dropout mechanism to its low-rank updates to regularize and avert overfitting. While the standard Dropout technique~\cite{srivastava2014dropout} readily applies to LoRA's additive weight modifications, it does not seamlessly adapt to the multiplicative nature of BOFT updates. To resolve this, we introduce a multiplicative Dropout scheme for BOFT. Since $\bm{R}(m,b)$ consists of $m$ orthogonal butterfly modules, each rearrangeable into $2b\times 2b$ block-diagonal orthogonal forms, our multiplicative Dropout method works by first selecting at random $p_1$ percent of the butterfly modules, and within each selected module, further picking $p_2$ percent of their diagonal blocks to be replaced by identity matrices.

\

\textbf{Expressivity of BOFT}. The combination of the butterfly matrix structure and permutations can exactly reconstruct various well-known fast linear transforms~\cite{dao2019learning,dao2019kaleidoscope} (\emph{e.g.}, fast Fourier transform, Hadamard transform). However, the extent to which our orthogonal butterfly matrix can serve as an approximation for a general orthogonal matrix remains to be fully established. To address this, we performed simulations aimed at approximating a randomly generated dense orthogonal matrix~\cite{anderson1987generation} of dimension $512\times 512$. Figure~\ref{fig:para_eff} displays averages from 10 runs with distinct random seeds. The vertical axis measures the error in approximation, while the horizontal axis reports the count of effective trainable parameters. Each identically colored curve corresponds to BOFT configurations with equal block size, where the initial point (at the far left) corresponds to BOFT($1$,$b$), i.e., the standard block-diagonal OFT using block size $b$. Our empirical results show that BOFT consistently achieves superior parameter efficiency relative to OFT. For instance, BOFT($9$,$2$) achieves lower error than BOFT($1$,$16$) despite utilizing significantly fewer parameters. Configurations with smaller block size $b$ and increased depth $m$ tend to provide improved parameter efficiency; BOFT($6$,$4$) demonstrates comparable approximation error to BOFT($2$,$16$) but at a reduced parameter cost. In general, the butterfly construction captures a richer, more constrained subset within the orthogonal group than block-diagonal structures, leading to BOFT being theoretically more expressive than OFT with equivalent block size.

\begin{theorem}[Expressivity of BOFT]\label{thm:exp_boft}
    BOFT surpasses OFT in expressivity when both use the same block size. To approximate any $d\times d$ orthogonal matrix, it is sufficient to use a product of butterfly matrices: $\bm{B}_{d-1,1}(d)\bm{B}^\top_{d-1,2}(d)\cdots\bm{B}_{1,1}(d)\bm{B}^\top_{1,2}(d)$, with each $\bm{B}_{i,j}(d)$ a butterfly matrix for all valid $i$ and $j$.
\end{theorem}

Theorem~\ref{thm:exp_boft} points to a straightforward extension for BOFT: the final orthogonal matrix may be formulated as $\thickmuskip=2mu \medmuskip=2mu \bm{R}^G(m_1,b_1,m_2,b_2,l)=\bm{R}_{l,1}(m_1,b_1)\bm{R}_{l,2}^T(m_2,b_2)\cdots\bm{R}_{1,1}(m_1,b_1)\bm{R}_{1,2}^T(m_2,b_2)$, where $\bm{R}_{i,j}^T(m,b)$ are the orthogonal matrices implemented within BOFT. Choosing $m_1=m_2=\log d$, $b_1=b_2=2$, and $l=d-1$ makes $\bm{R}^G(m_1,b_1,m_2,b_2,l)$ capable of parameterizing every element of the orthogonal group, an approach also referred to as the kaleidoscope hierarchy~\cite{dao2019kaleidoscope}. It is important, however, to recognize that increased expressivity does not invariably yield improved finetuning results; in practice, full finetuning—while capable of universal expressivity—can often underperform. The crucial factor in the success of finetuning methods lies in balancing the degree of expressivity against model regularity. By equipping the finetuning parameter space with structural priors, BOFT allows practitioners to navigate this balance more effectively, increasing the chances of striking an optimal trade-off between flexibility and regularity.

\textbf{Spectral properties}. Orthogonal finetuning typically demonstrates superior spectral characteristics compared to LoRA, primarily because it maintains the spectral norm of the pretrained weight matrix $\bm{W}^0$ exactly. This can be demonstrated using singular value decomposition: $\bm{W}^0=\bm{U}\bm{\Sigma}\bm{V}^\top$, where $\bm{U}$ and $\bm{V}$ are orthogonal matrices and $\bm{\Sigma}$ contains the singular values along its diagonal. In both OFT and BOFT, an orthogonal matrix $\bm{R}$ is multiplied from the left, leading to finetuned weights of the form $\bm{R}\bm{U}\bm{\Sigma}\bm{V}^\top$, which preserves the largest singular value (i.e., the spectral norm) of $\bm{W}^0$. This spectral norm preservation is known to facilitate both stable training dynamics and improved generalization~\cite{miyato2018spectral,yoshida2017spectral}. Additional mathematical insights on this property can be found in Appendix~\ref{app:sn_property}.

\textbf{Orthogonal finetuning as learning bilinear similarity}. The BOFT approach can be reformulated as learning a bilinear similarity $\bm{w}^0_i\bm{R}\bm{x}$, where $\bm{w}^0_i$ denotes the $i$-th column (or neuron) of the pretrained weight matrix $\bm{W}^0$. Under this framework, BOFT amounts to learning the bilinear similarity transformation matrix $\bm{R}$ under the strong constraint that $\bm{R}$ must remain orthogonal. This perspective aligns with classical distance metric learning~\cite{xing2002distance} as well as with the study of bilinear forms~\cite{roman2005advanced}.

\textbf{Inductive bias and generalization in BOFT}. In the BOFT paradigm, $\bm{R}(m,b)$ typically denotes a structured subclass of the orthogonal group, thus restricting the class of admissible hypotheses. This structure inherently introduces an inductive bias. The specific structured inductive bias arising from butterfly factorization is posited to be favorable for generalization, since it encapsulates architectural motifs prevalent in traditional linear operators~\cite{dao2019kaleidoscope}—including the discrete Fourier transform, discrete sine/cosine transform, and the Hadamard transform. Furthermore, the use of sparse matrix factorization in BOFT may confer additional, implicit inductive biases~\cite{gunasekar2017implicit,li2018algorithmic,liu2019neural}.

\textbf{Contrast with butterfly-parameterized sparse training}. A substantial body of literature~\cite{dao2019learning,dao2019kaleidoscope,chen2022pixelated,dao2022monarch} has investigated sparse neural network training via the butterfly parameterization. The majority of these studies emphasize directly replacing weight matrices with butterfly-structured ones and learning models from the ground up. In contrast, \cite{dao2022monarch} explores finetuning by projecting pretrained weights onto butterfly-structured matrices and then adapting the projected subspaces for specific tasks. BOFT introduces an alternative finetuning paradigm, wherein transformation of weights is achieved using layer-wise shared matrices.

\input{rebuttal-results}

\section{Concluding Remarks and Limitations}

In this work, we present Orthogonal Butterfly, a versatile and parameter-efficient finetuning technique for foundation models, leveraging butterfly-based matrix structures. Our central contribution lies in achieving higher parameter efficiency by expressing a dense orthogonal matrix as the product of multiple sparse orthogonal matrices. To systematically identify viable factorizations, we introduce a graphical information transmission perspective. Within this paradigm, we establish that butterfly structures are especially effective for factorizing sparse orthogonal matrices. Empirical evaluations reveal the robustness and utility of BOFT for finetuning large-scale language models, visual models, and text-to-image generators. Results consistently underscore BOFT's capability as a broad finetuning method.

Nonetheless, BOFT comes with certain limitations. Because BOFT represents the final orthogonal matrix as a product of several orthogonal factors, training is somewhat slower compared to OFT, introducing additional computational cost during learning. Addressing the efficiency of BOFT’s training procedure thus remains an area for future research. On the positive side, post-finetuning, the composed orthogonal matrices can be fused into the pretrained network, incurring no additional latency during inference. Additionally, the question of whether butterfly networks optimally transmit information remains unresolved. Our information transmission framework opens opportunities to draw on insights from fields such as computer networking, where the efficacy of different network topologies for information flow is well-studied. This cross-disciplinary approach may yield more efficient designs for constructing dense orthogonal matrices.